\newpage
\mychapter{V. CONCLUSION AND RECOMMENDATION}{V. CONCLUSION AND RECOMMENDATION}

\section{Conclusion}

\hs This thesis set out to make Cambodian law easier to reach, and it did so by designing,
building, and evaluating an AI-powered legal assistant based on an Agentic
Retrieval-Augmented Generation (RAG) architecture. The motivation was simple but pressing:
legal information in Cambodia is public in principle yet hard to use in practice, written
in a formal Khmer that most people cannot easily interpret. The system addresses this by
letting a user ask a question in everyday Khmer and receive a concise, accurate answer that
points directly to the article of law it came from.

\hs At its core, the system is a pipeline of five stages connected through a directed state
graph in LangGraph: intent classification, query rewriting, hybrid retrieval, grounding
verification, and a two-step answer generation. The hybrid retrieval is what makes the
system robust in a low-resource setting, since it pairs dense semantic search using FAISS
with sparse keyword matching using BM25 and fuses the two with Reciprocal Rank Fusion. The
two-step generation, which separates writing the Khmer answer from extracting its
citations, keeps the prose fluent while still tying every answer back to a real source.

\hs The system was tested on a manually built dataset of 30 questions across 10 legal
categories and all four indexed labour laws. In its final form it reached a Retrieval
Precision@5 of 76.7\%, a Citation Accuracy of 70.0\%, and a Grounding Rate of 76.7\%.
These figures show that an agentic, hybrid RAG approach can work for legal question
answering in Khmer, despite the language's limited NLP support. The path to those numbers
is itself part of the result: targeted fixes to the intent classifier, the grounding
agent, and the query rewriter lifted Precision@5 from a baseline of 66.7\% to 76.7\%,
which confirms that each stage of the pipeline pulls real weight rather than simply adding
complexity.

\hs The work did not stop at a prototype on a laptop. The system was built into a complete
full-stack web application and deployed to a public server. It runs as Docker containers on
an Amazon EC2 instance behind Nginx with HTTPS, and it offers a React chat interface with
live answer streaming, citation chips, and an integrated PDF viewer that opens the source
law at the exact article page. The deployed system also includes optional authentication
through Clerk: anonymous visitors may ask a limited number of questions before being
invited to sign in, and signed-in users keep their conversation history. Throughout, the
architecture was kept domain-agnostic, so new areas of Cambodian law can be added simply by
running their documents through the data pipeline, without touching the agent system.

\hs In short, this work shows that Agentic RAG is a practical and effective way to build
domain-specific legal question answering for a low-resource language, and it leaves behind
a working, deployed prototype that takes a real step toward making Cambodian legal
information accessible to the people it affects.

% ─────────────────────────────────────────────────────────────────────────────
\section{Recommendations}

\hs Several practical lessons emerged during the development and evaluation of the system.
The following recommendations are offered to practitioners and researchers building
similar retrieval-augmented systems, particularly in low-resource language settings.

\begin{itemize}

    \item \textbf{Use asymmetric task types for embedding.} When an embedding model
    supports it, the document side and the query side should be embedded with different
    task types. Using \texttt{RETRIEVAL\_DOCUMENT} at indexing time and
    \texttt{RETRIEVAL\_QUERY} at query time, as done here, aligns the two far better than
    using a single task type for both, and noticeably improves retrieval.

    \item \textbf{Separate answer generation from citation extraction.} Asking a single
    model call to produce fluent prose and well-formed structured citations at the same
    time degrades both. Splitting the work into two calls, one for the Khmer answer and one
    for the citations, gives better results on each.

    \item \textbf{Give the grounding check enough context.} Judging relevance from only the
    first 200 characters of a chunk is unreliable for articles that open with a long
    preamble, since the substance may not appear within that window. Using at least 500
    characters per chunk reduced false-negative grounding decisions.

    \item \textbf{Keep query rewriting focused on one law domain.} A rewriter that injects
    several law-domain names into a single query adds noise and can pull keyword retrieval
    toward the wrong law. It is better to infer the single most relevant domain and target
    that one, as the regression in Run 2 made clear.

    \item \textbf{Cap the retry loop explicitly.} Without a hard limit, a poor query can
    send the grounding-and-rewrite loop round indefinitely. A maximum of two retries,
    followed by a generate-anyway fallback, keeps response time predictable while still
    giving retrieval a second chance.

\end{itemize}

% ─────────────────────────────────────────────────────────────────────────────
\section{Future Works}

\hs The system in its current form is a solid proof of concept, but several clear
directions remain for extending and strengthening it.

\begin{itemize}

    \item \textbf{QA-Augmented Indexing.} The system's main weakness is the vocabulary
    mismatch in the Social Security Law, which produced 0\% retrieval there. A promising fix
    is QA-augmented indexing: using a language model to generate two or three natural
    questions for each article, and then embedding those questions rather than the raw legal
    text. This would close the vocabulary gap at index time, so that a casual question about
    social-security contributions or work-injury benefits can match an embedded question
    instead of formal legal prose.

    \item \textbf{Extension to All Cambodian Law Domains.} The current knowledge base covers
    four labour laws. Because the architecture is domain-agnostic, future work can extend it
    to civil, criminal, commercial, land, and family law simply by running the data pipeline
    over additional official PDFs, with no change to the agent system itself.

    \item \textbf{Real-Time Token Streaming.} The present implementation simulates streaming:
    the full pipeline finishes before any text is sent to the user. Switching to true token
    streaming, using the Gemini SDK's \texttt{generate\_content\_stream}, would cut perceived
    latency and make the interface feel more responsive.

    \item \textbf{Cross-Encoder Reranking.} After RRF fusion, a cross-encoder could rerank
    the top candidates by jointly encoding the query and each article. Cross-encoders give
    more accurate relevance scores than bi-encoder similarity and could lift retrieval
    quality on the harder categories, at some additional computational cost.

    \item \textbf{Richer Account Features.} Authentication is already in place through Clerk,
    with per-user history and an anonymous free-tier limit. Building on this, future work
    could add saved or bookmarked answers, organisation accounts for firms and NGOs, and
    role-based access so that different user groups see features suited to their needs.

    \item \textbf{Multilingual Answer Generation.} The system currently answers only in
    Khmer. Adding English-language answers would open it to international users, legal
    researchers, and NGO staff who work with Cambodian law but do not read Khmer.

    \item \textbf{Feedback and Continuous Improvement.} A simple thumbs-up or thumbs-down on
    each answer would let the system gather human preference data over time. That data could
    then be used to fine-tune the query rewriter and the answer generator for the Cambodian
    legal domain, steadily improving quality with use.

    \item \textbf{Mobile Application.} The current interface is built for desktop browsers. A
    dedicated mobile app, or a fully responsive mobile layout, would widen access
    considerably, especially for workers in rural areas who reach the internet mainly through
    a smartphone.

\end{itemize}
